\chapter{Contrast Agents}

\label{ch:ch06_contrast_agents}

Contrast agents are substances introduced into the body to enhance the visibility of structures or fluids on imaging. They do not treat disease; they make anatomy and pathology more visible so that the correct diagnosis can be made. This chapter covers the major classes of contrast agents, how they work, how they are administered, and how to manage adverse reactions.


\section*{1. What Contrast Feels Like and Why It Is Used}
If you are a patient about to receive contrast, here is what to expect:
\begin{itemize}
    \item \textbf{Intravenous (IV) contrast}: You will feel a brief, warm flush spreading through your body and a metallic taste in your mouth. This is normal and passes within 30--60 seconds. You may also feel the need to urinate; this is the kidneys filtering the contrast and is not dangerous.
    \item \textbf{Oral contrast (barium or iodinated)}: You drink a chalky or mildly flavoured liquid. It may cause bloating or mild nausea. You may need to drink it over 1--2 hours before the scan.
    \item \textbf{Rectal contrast}: Administered as an enema; similar sensation to a standard enema.
    \item \textbf{Intra-articular contrast}: Injected directly into a joint (e.g., MRI arthrogram); local anesthetic is used first to minimize discomfort.
\end{itemize}

Contrast is used to:
\begin{itemize}
    \item Highlight blood vessels (angiography).
    \item Distinguish normal from abnormal tissue in organs such as the liver, kidneys, and brain.
    \item Show whether a structure is fluid-filled, solid, or vascular.
    \item Assess organ function (e.g., renal excretion, blood-brain barrier integrity).
\end{itemize}


\section*{2. Iodinated Contrast Agents}
Iodinated contrast is the most commonly used class of IV contrast worldwide, primarily for CT and fluoroscopic angiography.

\subsection*{How They Work}
Iodine has a high atomic number (53), which attenuates X-rays much more effectively than soft tissue. By attaching iodine atoms to an organic carrier molecule, a water-soluble compound is created that remains in the extracellular space and is excreted by the kidneys.

\subsection*{Classification}

\begin{table}[h]
\centering
\caption{Iodinated contrast agent classification}
\begin{tabular}{|p{0.15\textwidth}|p{0.11\textwidth}|p{0.21\textwidth}|p{0.37\textwidth}|}
\hline
\textbf{Type} & \textbf{Osmolality} & \textbf{Example} & \textbf{Key features} \\
\hline
Ionic, high-osmolality (HOCM) & $\sim$1500--2000 mOsm/kg & Diatrizoate (Urografin) & Largely historical; higher reaction rate. \\
\hline
Ionic, low-osmolality (LOCM) & $\sim$600--800 mOsm/kg & Ioxaglate (Hexabrix) & Lower reaction rate than HOCM. \\
\hline
Non-ionic, iso-osmolality (IOCM) & $\sim$290 mOsm/kg & Iodixanol (Visipaque) & Osmolality similar to blood; no established clinically meaningful kidney-safety advantage over modern low-osmolality agents for routine IV CT. \\
\hline
Non-ionic, low-osmolality (LOCM) & $\sim$600--800 mOsm/kg & Iohexol (Omnipaque), Iopamidol (Isovue), Iopromide (Ultravist) & Current standard of care; lowest reaction rate. \\
\hline
\end{tabular}
\end{table}

\begin{pearlbox}[Shellfish and contrast]
A shellfish allergy does not specifically predict an iodinated contrast reaction. The historical rationale (both contain iodine) is a myth. Allergies to shellfish are related to tropomyosin, not iodine. Always ask about prior contrast reactions, but do not assume shellfish allergy is a risk factor for contrast reactions.
\end{pearlbox}

\subsection*{Routes of Administration}
\begin{itemize}
    \item \textbf{Intravenous}: Most common. Bolus injection via a peripheral IV cannula (typically 18--20G in the antecubital fossa). Power injection rates range from 1--5 mL/s depending on the study.
    \item \textbf{Oral}: Used for CT abdomen/pelvis to opacify the GI tract. Iodinated water-soluble agents are preferred over barium when there is a risk of perforation.
    \item \textbf{Rectal}: Used in selected CT protocols (e.g., CT colonography with stool tagging).
    \item \textbf{Intra-articular}: For CT or MRI arthrography.
    \item \textbf{Intrathecal}: For myelography (CT myelogram).
\end{itemize}


\section*{3. Gadolinium-Based Contrast Agents (GBCAs)}
GBCAs are used in MRI to highlight vascularity, blood-brain barrier breakdown, and tissue perfusion. Gadolinium is a paramagnetic metal that shortens T1 relaxation time, causing bright signal on T1-weighted images.

\subsection*{ACR Risk Groups for Gadolinium-Based Contrast Agents}

\begin{table}[h]
\centering
\caption{ACR groups used to describe NSF risk; the groups are not a simple structural classification. Verify the current agent-specific listing before clinical use.}
\begin{tabular}{|p{0.12\textwidth}|p{0.30\textwidth}|p{0.44\textwidth}|}
\hline
\textbf{ACR group} & \textbf{NSF risk category} & \textbf{Examples} \\
\hline
Group I & Highest established NSF association; includes linear agents & Gadodiamide (Omniscan), gadoversetamide (OptiMARK), gadopentetate dimeglumine (Magnevist) \\
\hline
Group II & Very low NSF risk; includes agents with different structures & Gadoterate (Dotarem), gadobutrol (Gadavist), gadoteridol (ProHance), gadobenate dimeglumine (MultiHance) \\
\hline
Group III & Limited data; agent-specific guidance applies & Gadoxetate disodium (Eovist/Primovist) \\
\hline
\end{tabular}
\end{table}

\begin{warnbox}[Nephrogenic Systemic Fibrosis (NSF)]
NSF has been strongly associated with ACR Group I agents in patients with severe renal impairment. Group II agents have a very low risk, and a clinically indicated Group II examination generally should not be withheld solely because eGFR is below 30 mL/min/1.73 m$^2$. Group III agents have more limited data and require agent-specific assessment. Use the current ACR group listing and local policy; do not infer risk from linear or macrocyclic structure alone.
\end{warnbox}

\subsection*{Gadolinium Deposition}
Gadolinium retention in tissues, including the brain, has been reported after GBCA administration and is generally greater with some linear agents than with macrocyclic agents. No adverse health effect has been established from retention in patients with normal renal function. Use contrast when clinically justified and avoid unnecessary repeated doses; do not delay needed imaging solely because of retention concerns.
\begin{itemize}
    \item Agent selection according to the indication, current safety guidance, and local formulary.
    \item Individual assessment when repeated contrast-enhanced examinations are anticipated.
    \item Recognition that regulatory restrictions differ by country and agent.
\end{itemize}


\section*{4. Microbubble Contrast (Ultrasound)}
Microbubble contrast agents (e.g., SonoVue/Lumason [sulphur hexafluoride], Definity [perflutren]) consist of tiny gas-filled microspheres (1--10 $\mu$m) stabilised by a phospholipid shell. They remain entirely within the intravascular space and oscillate when exposed to ultrasound, producing a strong non-linear signal that dramatically improves blood pool visualisation.

Clinical applications:
\begin{itemize}
    \item Characterisation of liver lesions (FLL-5 classification).
    \item Evaluation of cardiac wall motion and myocardial perfusion.
    \item Vascular assessment (e.g., testicular, carotid).
    \item Avoidance of CT/MRI in patients who cannot receive iodinated or gadolinium contrast.
\end{itemize}

Microbubble agents have an excellent safety profile. Severe adverse reactions are very rare.


\section*{5. Barium and GI Contrast}
Barium sulphate is a radiopaque suspension used to opacify the gastrointestinal tract for fluoroscopic and CT examinations. It is not absorbed systemically.

\begin{warnbox}[When not to use barium]
Barium should not be administered if there is a suspected gastrointestinal perforation, as it can cause severe chemical peritonitis. Use water-soluble iodinated contrast instead (e.g., Gastrografin). Barium should also be avoided before certain endoscopic procedures.
\end{warnbox}

\begin{table}[h]
\centering
\caption{Barium vs water-soluble GI contrast}
\begin{tabular}{|p{0.21\textwidth}|p{0.30\textwidth}|p{0.35\textwidth}|}
\hline
\textbf{Feature} & \textbf{Barium sulphate} & \textbf{Water-soluble iodinated} \\
\hline
Aspiration risk & Can cause granuloma/chemical pneumonitis & Less irritating but hypertonic (may cause pulmonary edema) \\
\hline
Image quality & Excellent mucosal coating & Inferior mucosal coating \\
\hline
Safety in perforation & Contraindicated (peritonitis) & Safe \\
\hline
Typical use & Upper GI series, barium enema, CT colonography & Suspected perforation, post-operative checks \\
\hline
\end{tabular}
\end{table}


\section*{6. Patient Safety: Preparation and Screening}

\subsection*{Kidney Function}
Kidney-risk assessment depends on the contrast agent, route, current kidney function, acute kidney injury (AKI), and clinical need. Follow current institutional screening policy:
\begin{itemize}
    \item \textbf{Intravenous iodinated contrast}: In patients with stable eGFR $\geq30$ mL/min/1.73 m$^2$ and no AKI, the risk attributable to contrast is low. For AKI or eGFR $<30$, discuss risks and benefits; isotonic IV fluid prophylaxis may be appropriate if not contraindicated. Do not reduce contrast below a dose needed for a diagnostic examination.
    \item \textbf{GBCAs}: Use the current ACR agent risk group. Routine renal-function screening may not be required for Group II agents, but local policy and individual circumstances apply. AKI or severe CKD warrants agent-specific risk-benefit assessment rather than automatic denial of clinically needed MRI.
\end{itemize}

\begin{warnbox}[Metformin and contrast]
For intravenous iodinated contrast, metformin generally need not be stopped when eGFR is $\geq30$ mL/min/1.73 m$^2$ and there is no AKI. With AKI, eGFR $<30$, or selected intra-arterial procedures with risk of renal embolization, withhold it at or before the procedure and reassess kidney function before restarting, following current guidance and local policy.
\end{warnbox}

\subsection*{Allergy and Prior Reaction Screening}
\begin{itemize}
    \item Ask about prior contrast reactions (type, severity, agent, date).
    \item Document prior reactions to the same contrast class, including symptoms, severity, timing, and the specific agent; this is the most useful allergy-related history.
    \item Other allergies and asthma do not by themselves contraindicate contrast or predict a severe reaction; assess them in context and follow local screening policy.
\end{itemize}

\subsection*{Pregnancy and Breastfeeding}
\begin{itemize}
    \item \textbf{Iodinated contrast}: Use during pregnancy when clinically needed and the expected diagnostic benefit justifies it; pregnancy alone is not a reason to withhold an indicated examination. Breastfeeding generally does not need to be interrupted after intravascular iodinated contrast.
    \item \textbf{GBCAs}: Avoid routine administration during pregnancy; consider only when essential to diagnosis and expected to improve maternal or fetal outcome, after individualized risk-benefit assessment.
    \item \textbf{Ultrasound contrast}: Pregnancy use is agent- and indication-specific; consult current product information and local guidance rather than assuming universal safety.
\end{itemize}


\section*{7. Adverse Reactions: Classification and Management}

\subsection*{ACR Contrast Reaction Classification}

\begin{table}[h]
\centering
\caption{ACR classification of contrast reactions}
\begin{tabular}{|p{0.18\textwidth}|p{0.37\textwidth}|p{0.35\textwidth}|}
\hline
\textbf{Severity} & \textbf{Examples} & \textbf{Management} \\
\hline
Death & Refractory cardiovascular collapse & CPR, advanced cardiac life support. \\
\hline
Mild (self-limited) & Nausea, vomiting, brief urticaria, mild flushing, pruritus, warmth at injection site, mild nasal congestion & Observation; treat symptoms if needed (e.g., antihistamines for urticaria). No specific treatment usually required. \\
\hline
Moderate & Prolonged urticaria ($>$30 min), facial/laryngeal edema without respiratory compromise, bronchospasm with SpO$_2$ $>$90\%, hypertension (SBP 80--100 mmHg) & IV fluids, diphenhydramine, corticosteroids; consider nebulized salbutamol for bronchospasm. Monitor closely for progression. \\
\hline
Severe (life-threatening) & Hypotension (SBP $<$80 mmHg), respiratory arrest, cardiac arrest, laryngeal edema with stridor, severe bronchospasm with SpO$_2$ $<$90\%, arrhythmias & Immediate emergency response: oxygen, IV fluids, IM adrenaline (epinephrine), advanced airway management, CPR as needed. Follow local anaphylaxis protocol. \\
\hline
\end{tabular}
\end{table}

\subsection*{Adrenaline (Epinephrine) Dosing for Severe Reactions}

\begin{table}[h]
\centering
\caption{Adrenaline dosing for contrast reactions}
\begin{tabular}{|p{0.35\textwidth}|p{0.25\textwidth}|p{0.26\textwidth}|}
\hline
\textbf{Presentation} & \textbf{Dose} & \textbf{Route} \\
\hline
Bronchospasm (no hypotension) & 0.3--0.5 mg (1:1000) & IM (anterolateral thigh) \\
\hline
Cardiac arrest & 1 mg (1:10,000) & IV (per ACLS protocol) \\
\hline
Hypotension (SBP $<$80 mmHg) & 0.1--0.3 mg (1:10,000) & IV slow push \\
\hline
\end{tabular}
\end{table}

\begin{pearlbox}[Acute contrast reactions]
Use the terms \textbf{allergic-like} and \textbf{physiologic} for acute iodinated-contrast reactions, as in current radiology guidance. An allergic-like reaction can occur without a known IgE mechanism, and clinical management is based on presentation and severity. A prior reaction to the same contrast class is the most important predictor of recurrence.
\end{pearlbox}

\subsection*{Prior Contrast Reactions}
For a patient with a prior contrast reaction, first clarify the reaction and consider an alternative examination or a different agent in consultation with radiology. Corticosteroid premedication is not routinely recommended after a prior mild immediate reaction; selected patients with more serious prior reactions may warrant a locally approved regimen when contrast is necessary. Premedication does not prevent all reactions and must not delay urgent imaging or treatment. If a facility uses a regimen, follow its current protocol and remain prepared to treat a breakthrough reaction.


\section*{8. Acute Kidney Injury After Contrast}

\begin{itemize}
    \item \textbf{Contrast-associated AKI (CA-AKI)} means AKI occurring after contrast exposure, whether or not contrast caused it. \textbf{Contrast-induced AKI (CI-AKI)} refers to injury causally attributable to contrast; it is difficult to distinguish from other causes in an individual patient.
    \item The risk of CI-AKI after intravenous iodinated contrast is low for most patients with stable eGFR $\geq30$ mL/min/1.73 m$^2$; uncertainty and risk are greatest with AKI or eGFR $<30$.
    \item For AKI or eGFR $<30$, consider isotonic IV saline prophylaxis if not contraindicated, and weigh diagnostic benefit against individual risk. Do not use an arbitrarily reduced contrast dose that could make the examination nondiagnostic.
    \item Low-osmolality and iso-osmolality agents have no established clinically meaningful difference in CI-AKI risk for routine intravenous CT use.
\end{itemize}

\begin{pearlbox}[CIN vs NSF: different risks, different agents]
CA-AKI and NSF are distinct. NSF is chiefly associated with particular GBCA risk groups in patients with severe renal impairment; it should not be generalized to every gadolinium agent or described as a risk from linear structure alone.
\end{pearlbox}


\section*{9. eGFR Thresholds and Decision-Making}

\begin{table}[h]
\centering
\caption{Simplified kidney-risk framework; use agent-specific current guidance and local policy.}
\begin{tabular}{|p{0.15\textwidth}|p{0.36\textwidth}|p{0.36\textwidth}|}
\hline
\textbf{eGFR} & \textbf{Iodinated contrast} & \textbf{GBCAs} \\
\hline
$\geq$30, stable and no AKI & Proceed when indicated; routine prophylaxis is generally unnecessary & Use an agent appropriate to the indication and current ACR group \\
\hline
$<$30 or AKI & Individual risk-benefit discussion; consider isotonic IV saline if not contraindicated; do not compromise diagnostic dose & A clinically indicated Group II agent should not automatically be withheld; assess agent, indication, and current guidance \\
\hline
Maintenance dialysis & Do not initiate or reschedule dialysis solely because intravenous iodinated contrast was given & Do not initiate dialysis or alter its schedule solely for a Group II GBCA; coordinate if otherwise clinically indicated \\
\hline
\end{tabular}
\end{table}


\section*{10. Emergency Equipment and Monitoring}
Facilities administering contrast should maintain emergency equipment, trained personnel, and a documented response pathway appropriate to the agents and procedures used, in accordance with current standards and local regulation. This typically includes:
\begin{itemize}
    \item Oxygen supply and delivery equipment (masks, nasal cannulae).
    \item Suction.
    \item Emergency medications: adrenaline (epinephrine), diphenhydramine, corticosteroids, salbutamol nebuliser, IV fluids.
    \item Blood pressure cuff and pulse oximeter.
    \item Resuscitation trolley.
    \item Access to emergency call system.
\end{itemize}

Staff must be trained in contrast reaction recognition and management, with regular drills. A documented reaction pathway should be available in every imaging department.

\begin{pearlbox}[The most dangerous reaction is the one you do not recognize]
Most severe immediate hypersensitivity reactions occur within the first 15--30 minutes, so staff and facilities should follow current monitoring and discharge policies, with heightened attention to patients with prior reactions. A brief warmth or metallic taste is a common physiologic effect, but new or worsening symptoms require assessment.
\end{pearlbox}

\section*{Selected Guidance}
\begin{itemize}
    \item American College of Radiology. \href{https://www.acr.org/Clinical-Resources/Clinical-Tools-and-Reference/Contrast-Manual}{Manual on Contrast Media} (current agent groups, kidney risk, pregnancy, and acute reactions).
    \item American College of Radiology and National Kidney Foundation. \href{https://doi.org/10.1148/radiol.2019192094}{Consensus on intravenous iodinated contrast in patients with kidney disease}.
    \item American College of Radiology and National Kidney Foundation. \href{https://doi.org/10.1148/radiol.2020202903}{Consensus on intravenous gadolinium-based contrast in patients with kidney disease}.
    \item American College of Radiology and American Academy of Allergy, Asthma \& Immunology. \href{https://doi.org/10.1148/radiol.240100}{Consensus on prevention of recurrent iodinated-contrast hypersensitivity reactions}.
\end{itemize}

\clearpage
